\chapter{Bending Spoons seminary (Evernote)}

Founded in 2013 in Danmark.
They acquire digital products then unlock the potential of the products.


500+ TB data processed per year. Constantly search for new talent.

They look for:

\begin{itemize}
  \item Drive
  \item Potential
  \item Passion
\end{itemize}

The processes:

\begin{itemize}
  \item Screening
  \item Tasks
  \item Interviews
  \item Final decision
\end{itemize}


\section{Evernote}

There is a frontend that get the data from bakend at the very high level, but deeper
the system is very complex and full of components.

Before the acquisition Evernote was a monolith, but after the acquisition
they decided to migrate into a set of microservices with each one
thake care of a specific element.

All of there service shoudl minimize the dependecies with the other services.

But the backend is not the complex part of the system.

\subsection{Evernote's sync system}

Legacy: "Monolith have you some new content?" "Yes, the note with the monkey"
This is a polling request useless for the user and for the backend system.


Microservices: Persistent connection between the client and the server, faster
then polling, cheaper for both:  the client side save battery and for
the server side save CPU and bandwidth.


\subsubsection{Deep dive into sync activity}

Note creation:


\begin{itemize}
  \item The phone perform an http request to the backend to note service
  \item Performing some checks
  \item Create the json file SQL database
  \item Note create message to the sync service queue
  \item The sync service send the message to the all the client devices with the new note
\end{itemize}



\subsubsection{Client connection}
Fernando autenticate withe the websocket and then connect to sycn service.
This service has tons of persistent connection always open, it use a
multi-map, because of multiple devices but same user, to keep track of
the connection and the user id.


The persistent connection can go offline for a bit. Do they loose all the
message? Yes but the sync service use google cloud datastore.
It is a single schema with the UId, the Instance with every metadata of
Fernando and SyncOperation with the operation performed (CREATE, UPDATE, DELETE).

Every time the client receaves some message, it reaceve the userId and
the timestamp and this is used in the google cloud datastore.

\begin{minted}[breaklines]{sql}
SELECT * FROM sync_documents
WHERE accessInfo.updated >= "1259414" and accessInfo.userId = "12345"

\end{minted}

The sync service have some pros:

\begin{itemize}
  \item Fast updates
  \item Offline access
  \item Cheap
\end{itemize}


But also some drawsbacks:

\begin{itemize}
  \item Data duplication
  \item Integration difficulty
  \item Complexity
\end{itemize}





